\section{Legation, an Unpromulgated Index, and the Vulgate}

\subsection{Paris under siege}

From 1588 to 1592 Bellarmine served as spiritual father at the Roman College,
though these years also carried him beyond Rome.  In 1589 he accompanied a
papal legation to France and experienced the siege of Paris.  His
autobiography gives the itinerary and reads survival and service through
providence.  That spiritual reading belongs to the subject's evidence; it is
not a neutral account of every political faction or the population's
experience.

The legation placed a theologian of authority inside an acute contest over
dynasty, confession, papal policy, and war.  Bellarmine's written distinctions
about ecclesial and temporal authority were not academic abstractions.  They
could shape counsel toward rulers and judgments about allegiance.  The French
experience also demonstrated the limits of argument where institutions
possessed armies, courts, and censorship.

\subsection{Sixtus V and a book not formally condemned}

In 1590 Bellarmine himself came under papal suspicion.  Sixtus~V objected to
the \emph{Disputationes}, particularly the denial of direct papal dominion in
temporal affairs, and moved to restrict the work until correction.  Bellarmine
later narrated the conflict in his autobiography.  His account is direct
evidence of remembered danger and an interested defense of his position.

Modern reconstruction finds that the intended Sixtine Index was not
promulgated before Sixtus died.  The procedural stage and exact text still
require fuller archival collation.  It is therefore inaccurate to say that a
published universal Index formally condemned Bellarmine; it is equally
inaccurate to erase a serious attempt at papal censorship because it did not
reach promulgation.  The episode exposes genuine disagreement within the
Roman institutions Bellarmine served.

\subsection{Correction and papal reputation}

The next textual crisis concerned the Latin Bible.  A Vulgate edition produced
under Sixtus required substantial correction.  Bellarmine took part in the
work leading to the Sixto-Clementine edition of 1592 and advised a preface that
protected Sixtus's reputation by assigning faults substantially to printing
and correction.  This account comes from Bellarmine's own autobiography and is
confirmed in broad sequence by critical biography.

The event should be described without either polemical extreme.  Bellarmine did
not simply forge Scripture, and the problem was not presented by him as a few
irrelevant slips.  He contributed to real textual correction while also
managing the public honor of a deceased pope.  The case shows a scholar and
institutional servant confronting a tension between accuracy and reputation.
Without full collation of the Sixtine and Sixto-Clementine printings, this
study does not assign every variant or commission decision to him.

These episodes place two kinds of authority side by side.  A theologian's book
could be threatened by the very papacy it defended, and a papal biblical
edition could require correction by subordinate scholars.  Bellarmine did not
respond by abandoning the institution.  He worked through revision,
distinction, and protective explanation---a pattern that displays loyalty and
skill while leaving the ethics of candor open to judgment.

\evidenceaxis{Sixtus V moved to restrict Bellarmine's work in 1590, and
Bellarmine later helped correct the Sixtine Vulgate.}{Bellarmine's
autobiography at fols.~10v--11v and named modern reconstruction control both
episodes.}{The Index project was unpromulgated; the exact procedure and full
Vulgate edition history remain uncollated, while Bellarmine's account is also
self-defense and reputation management.}
